% ==================== 1 引言 ====================
\section{引　言}

X 射线衍射（X-ray diffraction，XRD）是纳米材料物相鉴定和结构表征中最常用的方法之一\cite{liang1991xray,cullity2001elements}。当一束单色 X 射线照射到晶体上时，原子中的电子受迫振动，向各方向发出相干散射波；只有在满足布拉格方程
\begin{equation}
2d_{hkl}\sin\theta=n\lambda
\label{eq:bragg}
\end{equation}
的方向上，散射波才相长叠加，形成探测器可以记录到的衍射峰。式中 $d_{hkl}$ 为 $(hkl)$ 晶面的面间距，$\theta$ 为布拉格角，$\lambda$ 为入射波长，$n$ 为衍射级数。

一张粉末衍射图至少包含峰位、峰强、峰宽和背底四类信息。对本次作业来说，最重要的是峰位和相对强度：峰位由晶面间距决定，可以用来和标准 PDF 卡片对照；相对强度反映晶胞内原子的种类、位置和占有率，是物相鉴定的辅助依据。PDF 卡片是 JCPDS/ICDD 收集的标准粉末衍射数据，每张卡片给出某一物相的 $2\theta$、$d$ 值、相对强度和 $(hkl)$ 指数\cite{pdf711167,pdf751756}。把实验谱和 PDF 卡片画在同一张图上，就能直观看出哪些峰来自哪个物相、有没有杂相。

TiO$_2$ 的常见晶型有锐钛矿（anatase）、金红石（rutile）和板钛矿（brookite）。三者的 TiO$_6$ 八面体连接方式不同，晶胞参数和衍射花样也不同；在热处理过程中，锐钛矿会逐渐向金红石转变，衍射图上表现为锐钛矿峰减弱、金红石峰出现并增强。本文使用的 8 组样品编号为 200--900，对应不同的热处理温度；数据由第四组同学于 2026 年 9 月 29 日采集，按课程要求在全班共享。本文只对导出的 txt 做格式整理、绘图和物相描述，不改动原始文件。

本次作业要完成两件事：一是按热处理温度把 8 组 XRD 数据画成多曲线图，并在图上叠加锐钛矿和金红石的 PDF 卡片；二是结合自己的实验过程，总结 XRD 实验的操作要点、注意事项、存在问题及解决办法。因此，本文不展开更深入的定量分析，这部分留到后续课程作业中再做。全文安排如下：第 2 节说明样品、仪器和测试参数；第 3 节说明原始数据的格式问题、\texttt{make\_\allowbreak origin\_\allowbreak ready.py} 的转换方法和 Origin 多曲线绘制；第 4 节说明 PDF 卡片的检索、绘制和叠加，并给出最终的物相描述；第 5 节总结实验操作与常见问题；第 6 节为结论。
